\documentclass[11pt,letterpaper]{article}
\usepackage[margin=2cm]{geometry}
\usepackage{amsmath,amssymb,graphicx,xcolor,booktabs,tabularx,enumitem,titlesec,fancyhdr,array,longtable}
\usepackage[hidelinks]{hyperref}
\usepackage{helvet}\renewcommand{\familydefault}{\sfdefault}
\titleformat{\section}{\Large\bfseries\color{blue!50!black}}{\thesection}{0.6em}{}
\titleformat{\subsection}{\large\bfseries}{\thesubsection}{0.6em}{}
\setlength{\headheight}{14pt}
\pagestyle{fancy}\fancyhf{}\lhead{sPHENIX Self-Paced Course}\rhead{Syllabus}\cfoot{\thepage}
\newcommand{\infobox}[2]{\par\noindent\fcolorbox{blue!50!black}{blue!4}{\begin{minipage}{0.95\linewidth}\textbf{#1}\par\smallskip #2\end{minipage}}\par\medskip}
\setlist{itemsep=1pt,topsep=3pt}
\setlength{\emergencystretch}{2em}
\begin{document}
\begin{center}{\LARGE\bfseries sPHENIX Self-Paced Course}\\[4pt]{\Large Syllabus and Study Plan}\\[8pt]
Prepared for P.\,D.\ Lewis (Ohio University) \quad November 2026\\[2pt]
\small Built on the Summer-2026 sPHENIX School and Workfest (BNL, 27 July -- 14 August 2026), public sPHENIX papers and your own notes\end{center}
\section{Purpose}
This course sharpens your working knowledge of the whole sPHENIX experiment: the physics programme, the accelerator, every detector
subsystem, the trigger and data acquisition, the Fun4All software and simulation, reconstruction and calibration, and the analysis
workflow from good-run list to publication. You already know the EMCal $\pi^0$/$\eta$ corner of the experiment well; the goal is to
understand the parts around it well enough to read any sPHENIX analysis note, ask good questions in meetings, and make better decisions in
your own Run-24 $p$+$p$ $\pi^0$/$\eta$ cross-section analysis.
\infobox{How the course is built}{Ten module PDFs, each 14--25 pages. Every module follows the same pattern: learning objectives and
prerequisites; numbered steps, each with \emph{goal}, \emph{physics}, \emph{how it is done in sPHENIX}, \emph{checks} and \emph{pitfalls};
``Connect to your analysis'' boxes; suggested hands-on exercises on SDCC; self-test questions with answers; a glossary and references.
The school talks behind each module are cited as [S\#\#] (list in Section~\ref{sec:talks}).}
\section{How to use this course}
\begin{enumerate}
\item Work through one module per week (two weeks for Modules 6 and 10). Budget the hours in Table~\ref{tab:plan}.
\item Each module is self-contained: it builds the material from first principles up to current sPHENIX practice, so you do not need the slides. The school talks it cites as [S\#\#] are optional further reading (Indico event 32182, collaboration login).
\item Read the module once straight through, then a second time working the checks. The SDCC exercises are optional.
\item Answer the self-test questions \emph{before} reading the answers. If you miss more than three, reread that step.
\item Keep a running log (one page per module) of what you learned and what you want to ask; bring it to group meetings.
\item After Module 10, write the capstone proposal and discuss it with your advisor.
\end{enumerate}
\section{Schedule}
\begin{table}[h]\centering\small
\begin{tabularx}{\linewidth}{@{}llXr@{}}\toprule
Week & Module & Title & Hours\\\midrule
1 & 1 & The sPHENIX physics program: QGP, jets, heavy flavour, cold QCD and spin & 8\\
2 & 2 & RHIC and the sPHENIX detector: layout, global detectors, runs and centrality & 8\\
3 & 3 & Calorimetry in sPHENIX: EMCal and HCal hardware & 8\\
4 & 4 & Tracking detectors: MVTX, INTT, TPC and TPOT & 8\\
5 & 5 & Trigger, DAQ and streaming readout & 7\\
6--7 & 6 & Fun4All framework, software stack and simulation & 14\\
8 & 7 & Calorimeter reconstruction and calibration: towers, clusters, photons, jets & 10\\
9 & 8 & Tracking reconstruction: seeding, fitting, alignment and TPC calibration & 9\\
10 & 9 & From data to a result: the sPHENIX analysis workflow & 10\\
11--12 & 10 & Case studies from the 2026 school and machine learning in sPHENIX & 14\\
\midrule
 & & Total & 96\\\bottomrule\end{tabularx}
\caption{Suggested 12-week plan (about 8--10 hours per week). The order follows the school: physics, detectors, software, reconstruction, analysis.}\label{tab:plan}\end{table}
\infobox{Dependencies}{Modules 1--2 are prerequisites for everything. Modules 3--5 can be read in any order. Module 6 is needed before
7 and 8. Module 9 uses 2, 5, 6 and 7. Module 10 uses all modules.}
\section{Modules}
\subsection*{Module 1: The sPHENIX physics program: QGP, jets, heavy flavour, cold QCD and spin}
\addcontentsline{toc}{subsection}{Module 1}
\textbf{File:} \texttt{Module01\_Physics\_Program.pdf} \quad \textbf{Week:} 1 \quad \textbf{Time:} 8 h\\
\textbf{Scope.} Why sPHENIX exists; QGP and jet quenching theory basics; history from PHENIX to sPHENIX; the four topical groups (Bulk, Jet Structure, Heavy Flavor, Cold QCD) and their flagship measurements; spin physics with polarized protons; where neutral mesons fit.\\
\textbf{You will be able to:}
\begin{itemize}
\item Explain what the QGP is and why jets, heavy flavour and quarkonia probe it
\item Describe the four topical groups and their flagship measurements
\item Explain the cold-QCD and spin programme with polarized protons
\item Place neutral-meson measurements ($\pi^0$, $\eta$, $R_{AA}$, cross sections) in the programme
\end{itemize}
\textbf{School talks:} [S01], [S02], [S03], [S04], [S05], [S06], [S07], [S44], [S50].\\
\textbf{Checkpoint exercise:} Read the two first Au+Au papers; write a one-page map of which sPHENIX measurement constrains which QGP property.
\medskip
\subsection*{Module 2: RHIC and the sPHENIX detector: layout, global detectors, runs and centrality}
\addcontentsline{toc}{subsection}{Module 2}
\textbf{File:} \texttt{Module02\_RHIC\_and\_Detector\_Overview.pdf} \quad \textbf{Week:} 2 \quad \textbf{Time:} 8 h\\
\textbf{Scope.} RHIC accelerator and beam basics (bunches, crossing angle, luminosity, polarization); detector layout and coordinates; solenoid; global detectors MBD, ZDC, sEPD; vertexing; centrality and Glauber; event plane; collision systems and runs 2023-2025.\\
\textbf{You will be able to:}
\begin{itemize}
\item Describe RHIC beams: bunches, crossing angle, luminosity, polarization
\item Draw the sPHENIX layout and coordinate system
\item Explain MBD, ZDC and sEPD and how vertex, centrality and event plane are obtained
\item Summarise the runs 2023--2025 and collision systems
\end{itemize}
\textbf{School talks:} [S01], [S08], [S48], [S49], [S05].\\
\textbf{Checkpoint exercise:} On SDCC, query the spin DB and DAQ DB for a run of your choice (crossing angle, scaledowns); plot the MBD vertex distribution for that run.
\medskip
\subsection*{Module 3: Calorimetry in sPHENIX: EMCal and HCal hardware}
\addcontentsline{toc}{subsection}{Module 3}
\textbf{File:} \texttt{Module03\_Calorimeters.pdf} \quad \textbf{Week:} 3 \quad \textbf{Time:} 8 h\\
\textbf{Scope.} Electromagnetic and hadronic showers; sampling calorimetry; EMCal tungsten-powder, scintillating-fiber blocks and SiPM readout; inner/outer HCal; towers, light guides, digitizers; energy resolution; test beams; calibration strategy overview (MIP, $\pi^0$, tower slope, cosmics).\\
\textbf{You will be able to:}
\begin{itemize}
\item Explain electromagnetic and hadronic shower development
\item Describe the EMCal and HCal construction and SiPM readout
\item Relate sampling fraction and stochastic/constant terms to resolution
\item Outline the calibration chain from test beam to in-situ
\end{itemize}
\textbf{School talks:} [S09], [S08], [S47], [S42].\\
\textbf{Checkpoint exercise:} Draw the EMCal tower grid; estimate the shower containment for a 2 GeV photon; look up one tower's calibration in the CDB.
\medskip
\subsection*{Module 4: Tracking detectors: MVTX, INTT, TPC and TPOT}
\addcontentsline{toc}{subsection}{Module 4}
\textbf{File:} \texttt{Module04\_Tracking\_Detectors.pdf} \quad \textbf{Week:} 4 \quad \textbf{Time:} 8 h\\
\textbf{Scope.} Charged-particle detection in silicon and gas; MVTX MAPS pixels; INTT strips and timing; TPC drift, gain (GEM), space charge; TPOT micromegas; integration times and pile-up; role of each in vertexing, momentum, timing.\\
\textbf{You will be able to:}
\begin{itemize}
\item Explain charge collection in silicon and gas detectors
\item Describe MVTX, INTT, TPC and TPOT geometry and roles
\item Explain integration time, timing and pile-up for each detector
\item Relate detector design to momentum and DCA resolution
\end{itemize}
\textbf{School talks:} [S10], [S11], [S12], [S28].\\
\textbf{Checkpoint exercise:} Sketch the hits of a 1 GeV pion through MVTX, INTT and TPC; estimate the TPC drift time and what it implies for pile-up.
\medskip
\subsection*{Module 5: Trigger, DAQ and streaming readout}
\addcontentsline{toc}{subsection}{Module 5}
\textbf{File:} \texttt{Module05\_Trigger\_DAQ\_Streaming.pdf} \quad \textbf{Week:} 5 \quad \textbf{Time:} 7 h\\
\textbf{Scope.} GL1 and trigger bits, MBD and calorimeter (photon/jet) triggers, scaledowns/prescales and live time; DAQ architecture and event building; triggered vs streaming readout; from raw data (PRDF) to DSTs; production, tags and the FileCatalog; SDCC computing access (ssh, web).\\
\textbf{You will be able to:}
\begin{itemize}
\item Explain GL1 trigger bits, scaledowns and live time
\item Describe MBD and calorimeter-based triggers
\item Contrast triggered and streaming readout
\item Follow data from PRDF to DSTs and find files in the FileCatalog
\end{itemize}
\textbf{School talks:} [S13], [S21].\\
\textbf{Checkpoint exercise:} Use the FileCatalog to list the DST segments of one run; read the GL1 scalers for that run.
\medskip
\subsection*{Module 6: Fun4All framework, software stack and simulation}
\addcontentsline{toc}{subsection}{Module 6}
\textbf{File:} \texttt{Module06\_Fun4All\_and\_Simulation.pdf} \quad \textbf{Week:} 6--7 \quad \textbf{Time:} 14 h\\
\textbf{Scope.} Fun4AllServer, SubsysReco lifecycle, node tree, input/output managers, macros, CDB and global tags, builds and installs (MYINSTALL), coresoftware/macros/analysis repos; Geant4 detector simulation (PHG4Reco), event generators (PYTHIA8, HIJING), single-particle guns and embedding, simulation productions; hands-on exercises.\\
\textbf{You will be able to:}
\begin{itemize}
\item Write and run a Fun4All macro with a custom SubsysReco
\item Navigate the node tree and the CDB/global tags
\item Build and install a module into a private area
\item Run Geant4 simulation with generators, particle guns and embedding
\end{itemize}
\textbf{School talks:} [S14], [S13], [S47].\\
\textbf{Checkpoint exercise:} Write a minimal SubsysReco that prints cluster counts; run a single-$\pi^0$ gun through Geant4 and reconstruct it.
\medskip
\subsection*{Module 7: Calorimeter reconstruction and calibration: towers, clusters, photons, jets}
\addcontentsline{toc}{subsection}{Module 7}
\textbf{File:} \texttt{Module07\_Calorimeter\_Reconstruction.pdf} \quad \textbf{Week:} 8 \quad \textbf{Time:} 10 h\\
\textbf{Scope.} Waveform fitting and TowerInfo; status and hot-tower maps; EMCal $\pi^0$ and HCal calibrations; clustering (template clustering, chi2/probability), topo clustering; photon identification and isolation; jet reconstruction (anti-kt, constituents, background subtraction, JES/JER).\\
\textbf{You will be able to:}
\begin{itemize}
\item Explain waveform fitting, TowerInfo and status maps
\item Describe EMCal and HCal calibration methods
\item Compare clustering algorithms and topo clusters
\item Explain photon ID/isolation and jet reconstruction with background subtraction
\end{itemize}
\textbf{School talks:} [S15], [S16], [S24], [S17], [S18], [S42].\\
\textbf{Checkpoint exercise:} Run clustering on one DST segment with two cluster settings and compare cluster counts and energies (no physics claims).
\medskip
\subsection*{Module 8: Tracking reconstruction: seeding, fitting, alignment and TPC calibration}
\addcontentsline{toc}{subsection}{Module 8}
\textbf{File:} \texttt{Module08\_Tracking\_Reconstruction.pdf} \quad \textbf{Week:} 9 \quad \textbf{Time:} 9 h\\
\textbf{Scope.} Clustering in silicon and TPC; seeding; ACTS Kalman fitting; vertexing; alignment (millepede-style); TPC distortions and calibration (lasers, space charge); track-calorimeter matching; performance evaluation (efficiency, momentum resolution, DCA).\\
\textbf{You will be able to:}
\begin{itemize}
\item Explain seeding and Kalman-filter fitting in ACTS
\item Explain detector alignment
\item Describe TPC distortions and their calibration
\item Explain track--calorimeter matching and performance metrics
\end{itemize}
\textbf{School talks:} [S19], [S20], [S22], [S23], [S36], [S27].\\
\textbf{Checkpoint exercise:} Run the tracking QA macros on a small sample and identify each stage's output nodes.
\medskip
\subsection*{Module 9: From data to a result: the sPHENIX analysis workflow}
\addcontentsline{toc}{subsection}{Module 9}
\textbf{File:} \texttt{Module09\_Analysis\_Workflow.pdf} \quad \textbf{Week:} 10 \quad \textbf{Time:} 10 h\\
\textbf{Scope.} Run QA and good-run lists; trigger selection and efficiency; luminosity and Vernier scans; centrality in practice; efficiency, acceptance and embedding; unfolding basics; systematic uncertainties (Barlow test); blinding/approval process (analysis notes, preliminaries, papers).\\
\textbf{You will be able to:}
\begin{itemize}
\item Build good-run lists and run QA
\item Compute luminosity and understand Vernier scans
\item Design efficiency/acceptance corrections and unfolding
\item Estimate systematic uncertainties and navigate the approval process
\end{itemize}
\textbf{School talks:} [S43], [S33], [S41], [S26], [S25], [S05], [S34].\\
\textbf{Checkpoint exercise:} Write the full error-budget table for your $\pi^0$/$\eta$ cross section and the run-selection flow chart.
\medskip
\subsection*{Module 10: Case studies from the 2026 school and machine learning in sPHENIX}
\addcontentsline{toc}{subsection}{Module 10}
\textbf{File:} \texttt{Module10\_Case\_Studies\_and\_ML.pdf} \quad \textbf{Week:} 11--12 \quad \textbf{Time:} 14 h\\
\textbf{Scope.} Worked mini-case-studies from student flash talks (multiplicity, flow $v_2$, jets, isolated photons and spin asymmetries, quarkonia energy correlators, phi-hadron correlations, CME, UPC rho, Vernier scans, INTT efficiency, HCal calibration); ML in physics (classification, regression for momentum scale, NN photon ID); capstone project plan tied to his $\pi^0$/eta analysis.\\
\textbf{You will be able to:}
\begin{itemize}
\item Analyse the structure of real student analyses from the school
\item Use ML for classification and regression in physics
\item Plan a capstone project tied to your $\pi^0$/$\eta$ analysis
\end{itemize}
\textbf{School talks:} [S30], [S31], [S32], [S34], [S35], [S40], [S37], [S38], [S39], [S29], [S27], [S25], [S26], [S41].\\
\textbf{Checkpoint exercise:} Write a two-page capstone proposal and present it to your advisor.
\medskip
\section{Capstone}
Write a two-page proposal for one improvement to your $\pi^0$/$\eta$ cross-section analysis that uses what you learned outside the EMCal
(for example: a track-based isolation or conversion veto, an MBD-trigger bias study with the global detectors, a streaming-readout
cross-check of luminosity, or an ML photon/$\pi^0$ discriminator for merged clusters). State the physics motivation, the data and
simulation needed, the method, the expected gain, and how you would validate it.
\section{Public readings}
\begin{itemize}
\item sPHENIX Public Results: \url{https://www.sphenix.bnl.gov/PublicResults}
\item sPHENIX Collab., charged hadron multiplicity in Au+Au, arXiv:2504.02240, JHEP 08 (2025) 075.
\item sPHENIX Collab., transverse energy density in Au+Au, arXiv:2504.02242, PRC 112 (2025) 024908.
\item sPHENIX Collab., dijet imbalance and acoplanarity in $p$+$p$, arXiv:2606.17184.
\item sPHENIX Collab., isolated prompt photons in $p$+$p$, arXiv:2607.03892.
\item INTT ladder and readout, NIM-A 1082 (2026) 171020, arXiv:2503.09105; INTT test beam, NIM-A 1086 (2026) 171312, arXiv:2509.00908.
\item Aidala et al., sPHENIX EMCal and HCal prototypes, IEEE Trans.\ Nucl.\ Sci.\ (2018) 2901--2919, arXiv:1704.01461.
\item PHENIX Collab., An Upgrade Proposal from the PHENIX Collaboration, arXiv:1501.06197.
\end{itemize}
\section{School talks cited in this course}\label{sec:talks}
All talks are from the Summer-2026 sPHENIX School and Workfest, BNL (Indico event 32182; collaboration access). The modules summarise
them in their own words; slides, plots and internal numbers are not reproduced.
\begin{longtable}{@{}lp{9.3cm}p{4.2cm}@{}}\toprule Key & Title & Speaker(s)\\\midrule\endhead
{[S01]} & Welcome to sPHENIX school and sPHENIX Overview & J. Huang, M. Connors\\
{[S02]} & sPHENIX History & D. Morrison\\
{[S03]} & sPHENIX Theory Overview: exploring QGP dynamics with jets & Y. Mehtar-Tani\\
{[S04]} & Cold QCD TG Analysis Opportunities & G. Nukazuka, M. Sarsour\\
{[S05]} & Bulk TG Analysis Opportunities & M. Mitrankova, R. Belmont\\
{[S06]} & Heavy Flavor TG Analysis Opportunities & A. Frawley, C. Dean, G. Ottino, M. Rosati\\
{[S07]} & Jet Structure TG Analysis Opportunities & J. Putschke, Y. Go\\
{[S08]} & HCal and Global Detector Overview & J. Haggerty\\
{[S09]} & EMCal Overview & C. Woody\\
{[S10]} & MVTX Overview & C. Dean\\
{[S11]} & INTT Overview & A. Enokizono\\
{[S12]} & TPC and TPOT Overview & T. Hemmick\\
{[S13]} & DAQ Overview (with web-browser and ssh tutorials) & S. Liechty\\
{[S14]} & Fun4All Overview / sPHENIX computing tutorial & C. Pinkenburg\\
{[S15]} & EMCal Reco and Calibrations & B. Seidlitz\\
{[S16]} & HCal Reco and Calibrations & A. Menon\\
{[S17]} & Photon Reconstruction & Y. Go\\
{[S18]} & Jet Reconstruction & V. Bailey\\
{[S19]} & Tracking Reconstruction & J. Osborn\\
{[S20]} & Tracking Alignment & G. Ottino\\
{[S21]} & Streaming Analysis & M. Peters\\
{[S22]} & TPC Calibration & M. Mitrankova\\
{[S23]} & Track-Calo Matching & X. Yu\\
{[S24]} & Topo Clustering & A. C. Oliveira da Silva\\
{[S25]} & Charged hadron multiplicity in O+O using silicon seeds (flash) & H.-R. Jheng\\
{[S26]} & Charged-hadron multiplicity in O+O collisions (flash) & C.-W. Shih\\
{[S27]} & Machine Learning Momentum Scale Correction (flash) & A. Patton\\
{[S28]} & INTT hit detection efficiency study (flash) & T. Kikuchi\\
{[S29]} & $J/\psi$ energy correlators in $p$+$p$ (flash) & B. Sayki\\
{[S30]} & Transverse single-spin asymmetries of isolated prompt photons (flash) & J. Hwang\\
{[S31]} & Neural network discrimination of isolated photons (flash) & R. Koike\\
{[S32]} & What I have done and what I plan to do next (flash) & S.-Y. Chen\\
{[S33]} & Run QA for Silicon-Calo Tracking Analysis (flash) & M. Shigihara\\
{[S34]} & Measurement of $v_2$ in Au+Au (flash) & I. Omae\\
{[S35]} & First Step in Jet Analysis (flash) & R. Takaaze\\
{[S36]} & Tracking Performance Evaluation in sPHENIX (flash) & A. P. Dash\\
{[S37]} & Constraining the CME using SDM (flash) & J. Singh\\
{[S38]} & Angular correlations of $\phi$ mesons and high-pT hadrons (flash) & K. Gable\\
{[S39]} & UPC $\rho^0$ production in O+O and Au+Au, simulation (flash) & L. Mwibanda\\
{[S40]} & Joint jet mass and charge for quark/gluon jet separation (flash) & M. I. Abdulhamid\\
{[S41]} & Vernier Scan analysis of $p$+$p$ and O+O in Run 25 (flash) & T. Harada\\
{[S42]} & HCal Calibration studies (flash) & A. M. Khan\\
{[S43]} & Jet Analysis & J. Park\\
{[S44]} & Cold QCD Lecture Series announcement & G. Nukazuka\\
{[S45]} & Closing & B. Sayki, J. Huang, M. Connors\\
{[S46]} & Workfest kickoff & S. Bathe, Y. Go\\
{[S47]} & Calorimeter test-beam simulation (workfest) & B. Seidlitz, J. Nagle, S. Bathe\\
{[S48]} & Spin dynamics at the RHIC accelerator complex (workfest lecture) & K. Hock\\
{[S49]} & Polarimetry at RHIC (workfest lecture) & X. Chu\\
{[S50]} & Cold-QCD lecture (workfest) & M. Cerutti\\
\bottomrule\end{longtable}
\end{document}
